\documentclass[notes,10pt, aspectratio=169]{beamer}

\usepackage{pgfpages}
% These slides also contain speaker notes. You can print just the slides,
% just the notes, or both, depending on the setting below. Comment out the want
% you want.
% Note visibility is selected after package loading at the end of this preamble.
%\setbeameroption{show only notes} % Only notes
%\setbeameroption{show notes on second screen=right} % Both

\usepackage{helvet}
\usepackage[default]{lato}
\usepackage{array}
\usepackage{eurosym}
\usepackage{setspace}
\usepackage{animate}
\usepackage{multicol}
\usepackage{breqn}
\usepackage{subcaption}
\usepackage{caption}
\captionsetup{font=scriptsize, justification=raggedright,singlelinecheck=false}
\usepackage[export]{adjustbox}
\usepackage[scale=2]{ccicons}
\usepackage{tikz}
\newenvironment{wideitemize}{\itemize\addtolength{\itemsep}{10pt}}{\enditemize}
\usepackage{natbib}

\usepackage{tikz}
\usepackage{verbatim}
\setbeamertemplate{note page}{\pagecolor{yellow!5}\insertnote}
\usepackage{amsmath}
\usepackage{bm}
%\usepackage{mathpazo}
\usepackage{hyperref}
\usepackage{lipsum}
\usepackage{multimedia}
\usepackage{multirow}
\usepackage{graphicx}
\usepackage{dcolumn}
\usepackage{bbm}
\newcolumntype{d}[0]{D{.}{.}{5}}

\usepackage{changepage}
\usepackage{appendixnumberbeamer}
% \newcommand{\beginbackup}{
%   \newcounter{framenumbervorappendix}
%   \setcounter{framenumbervorappendix}{\value{framenumber}}
%   \setbeamertemplate{footline}
%   {
%      \leavevmode%
%      \hline
%      box{%
%       \begin{beamercolorbox}[wd=\paperwidth,ht=2.25ex,dp=1ex,right]{footlinecolor}%
% %         \insertframenumber  \hspace*{2ex}
%       \end{beamercolorbox}}%
%      \vskip0pt%
%   }
%  }

 \setbeamertemplate{footline}{
    \hbox{%
    \begin{beamercolorbox}[wd=\paperwidth,ht=3ex,dp=1.5ex,leftskip=2ex,rightskip=2ex]{page footer}%
        %\usebeamerfont{title in head/foot}%
        %\insertshorttitle \hfill
        %\insertsection \hfill
        %\tableofcontents[currentsection] \hfill
        \insertsectionnavigationhorizontal{0.7\paperwidth}{}{} \hfill
        \insertframenumber{} / \inserttotalframenumber
    \end{beamercolorbox}}%
}
\newcommand{\backupend}{
   \addtocounter{framenumbervorappendix}{-\value{framenumber}}
   \addtocounter{framenumber}{\value{framenumbervorappendix}}
}


\usepackage{graphicx}
\usepackage[space]{grffile}
\usepackage{booktabs}

% These are my colors -- there are many like them, but these ones are mine.
\definecolor{blue}{RGB}{0,114,178}
\definecolor{red}{RGB}{213,94,0}
\definecolor{yellow}{RGB}{240,228,66}
\definecolor{green}{RGB}{0,158,115}

\setbeamercolor{section in head/foot}{fg=blue}

\useinnertheme[shadow]{rounded}
\setbeamercolor{block title example}{fg=black,bg=orange!50!white}
\setbeamercolor{block body example}{fg=black, bg=orange!30!white}


\hypersetup{
  colorlinks=false,
  linkbordercolor = {white},
  linkcolor = {blue}
}


%% I use a beige off white for my background
\definecolor{MyBackground}{RGB}{255,253,218}

%% Uncomment this if you want to change the background color to something else
%\setbeamercolor{background canvas}{bg=MyBackground}

%% Change the bg color to adjust your transition slide background color!
\newenvironment{transitionframe}{
  \setbeamercolor{background canvas}{bg=yellow}
  \begin{frame}}{
    \end{frame}
}

\setbeamercolor{frametitle}{fg=blue}
\setbeamercolor{title}{fg=black}
%\setbeamertemplate{footline}[frame number]
\setbeamertemplate{navigation symbols}{}
\setbeamertemplate{itemize items}{-}
\setbeamercolor{itemize item}{fg=blue}
\setbeamercolor{itemize subitem}{fg=blue}
\setbeamercolor{enumerate item}{fg=blue}
\setbeamercolor{enumerate subitem}{fg=blue}
\setbeamercolor{button}{bg=white,fg=blue,}



% If you like road maps, rather than having clutter at the top, have a roadmap show up at the end of each section
% (and after your introduction)
% Uncomment this is if you want the roadmap!
% \AtBeginSection[]
% {
%    \begin{frame}
%        \frametitle{Roadmap of Talk}
%        \tableofcontents[currentsection]
%    \end{frame}
% }
\setbeamercolor{section in toc}{fg=blue}
\setbeamercolor{subsection in toc}{fg=red}
%\setbeamersize{text margin left=1em,text margin right=1em}
\setbeamersize{text margin left=0.7cm,text margin right=0.7cm}


\title[]{\textcolor{blue}{Seminar 7:\\Functional Forms}}
\subtitle{Three problems and worked solutions}
\author{25117 - Econometrics}
\institute[shortinst]{Universitat Pompeu Fabra}

%\date{September 25th, 2024}


% Each problem is followed by worked solution slides. Set \showsolutionsfalse
% below to compile a problem-only deck. All data are hypothetical;
% no external assets are needed. Supplied uncertainty estimates support
% large-sample inference; no covariance-estimation formula is required.
\newif\ifshowsolutions
\showsolutionstrue
\newcommand{\E}{\mathbb{E}}
\newcommand{\Var}{\operatorname{Var}}
\newcommand{\Cov}{\operatorname{Cov}}
\newcommand{\Corr}{\operatorname{Corr}}
\newcommand{\answerpart}[1]{\textcolor{blue}{\textbf{#1.}}\,}

\setbeameroption{hide notes} % Speaker notes hidden; worked solutions are visible.


\date{}
\begin{document}
\begin{frame}[plain]
\maketitle
\end{frame}

\section{Polynomials}
\begin{frame}{Problem 1: Does an extra thousand euros have a constant association?}
\small
A hypothetical sample of 500 districts gives the quadratic fitted regression
\[
\widehat{Score}(x)=600+4x-0.04x^2,
\]
where $x$ is income in thousands of euros and scores are in points. Observed incomes range from 5 to 45. The supplied standard error of the coefficient on $X^2$ is $0.01$.
\begin{enumerate}
\item Calculate predicted scores at $x=10$ and $x=30$. Find the exact predicted score changes when income rises from 10 to 11 and from 30 to 31.
\item Derive the marginal effect $d\widehat{Score}(x)/dx$ and evaluate it at 10 and 30. Find the turning point. Should we use it to claim that scores decline at very high incomes?
\item Explain why OLS can estimate this model even though the relationship is curved. Test $H_0:\beta_{X^2}=0$ against a two-sided alternative at 5\%, using the large-sample normal approximation. What is the equivalent $F$ statistic for this one restriction?
\end{enumerate}
\end{frame}

\ifshowsolutions
\begin{frame}{Solution 1: Predictions and changing slopes}
\small
\answerpart{1}The fitted levels are
\[
\widehat{Score}(10)=600+40-4=636,\text{ and }
\widehat{Score}(30)=600+120-36=684.
\]
For any starting income $x$, a one-unit increase changes the prediction by
\[
\widehat{Score}(x+1)-\widehat{Score}(x)
=4-0.04[(x+1)^2-x^2]=4-0.04(2x+1).
\]
Thus the exact changes are $3.16$ points from 10 to 11 and $1.56$ points from 30 to 31.
\par\medskip
\answerpart{2}The marginal effect is
\[
\frac{d\widehat{Score}(x)}{dx}=4-0.08x.
\]
It is $3.20$ at $x=10$ and $1.60$ at $x=30$, in points per thousand euros. These derivatives describe an infinitesimal change; the exact one-unit changes average the declining slope over each interval.
\end{frame}

\begin{frame}{Solution 1: Turning point and the test of curvature}
\small
\answerpart{2}continued. Setting the marginal effect to zero gives
\[
4-0.08x=0\ \Longrightarrow\ x=50.
\]
The negative quadratic coefficient makes this a fitted maximum. However, 50 lies outside the observed income range $[5,45]$. A predicted decline beyond 50 is extrapolation and does not establish a decline in the population relationship.
\par\medskip
\answerpart{3}The model is linear in its coefficients:
\[
Score_i=\beta_0+\beta_1X_i+\beta_2X_i^2+U_i.
\]
Treat $X$ and $X^2$ as two regressor columns and estimate by OLS.
\[
t=\frac{-0.04}{0.01}=-4,\text{ and }F=t^2=16.
\]
Reject the zero-quadratic-coefficient null: $|t|>1.96$, or equivalently $F>3.84$ for one restriction. This is evidence of curvature within this specification.
\end{frame}
\fi

\section{Logarithms}
\begin{frame}{Problem 2: Which changes are in points and which are percentages?}
\small
Consider three hypothetical fitted functions for positive income $x$ (thousands of euros) and positive test score $y$. All logarithms are natural.
\[
\begin{aligned}
\text{A: }&\hat y=500+20\log x,\\
\text{B: }&\widehat{\log y}=6+0.02x,\\
\text{C: }&\widehat{\log y}=6+0.30\log x.
\end{aligned}
\]
\begin{enumerate}
\item Name each functional form and interpret its slope coefficient. Specify whether the change is in score points or percent.
\item For income rising from 10 to 11, calculate the exact change implied by each fitted function and the corresponding small-change approximation.
\item Repeat for a doubling from 10 to 20. Explain why some percentage approximations are less accurate for this larger change.
\end{enumerate}
For B and C, compare $\tilde y=\exp(\widehat{\log y})$. Use
$100[\exp(\Delta\widehat{\log y})-1]$ for its exact percentage change.
\end{frame}

\ifshowsolutions
\begin{frame}{Solution 2: Model A, step by step}
\small
\answerpart{1}A is linear--log: 1\% more income is associated with approximately $20/100=0.20$ more score points.
\par\smallskip
\answerpart{2}From 10 to 11, $\Delta x=1$ and $\Delta x/x=0.10$.
\begin{enumerate}
\item Evaluate $\hat y(x)=500+20\log x$ at the two incomes:
\[
\hat y(10)=500+20\log10,\text{ and }\hat y(11)=500+20\log11.
\]
\item Subtract. The intercept cancels, and $\log a-\log b=\log(a/b)$:
\[
\begin{aligned}
\hat y(11)-\hat y(10)&=20(\log11-\log10)\\
&=20\log(1.10)=\boxed{1.9062\text{ points}}.
\end{aligned}
\]
\item Approximate using $\log(1+g)\approx g$ for a small proportional change $g$:
\[
\Delta\hat y\approx20\frac{\Delta x}{x}=20(0.10)=\boxed{2\text{ points}}.
\]
\end{enumerate}
The outcome is in levels, so this is a change in score points.
\end{frame}

\begin{frame}{Solution 2: Model B, step by step}
\small
\setlength{\abovedisplayskip}{4pt}
\setlength{\belowdisplayskip}{4pt}
\answerpart{1}B is log--linear: one extra thousand euros is associated with approximately $100(0.02)=2$\% more in $\tilde y$.
\par\smallskip
\answerpart{2}Use $\widehat{\log y}(x)=6+0.02x$ and $\tilde y(x)=\exp(\widehat{\log y}(x))$.
\begin{enumerate}
\item Calculate the change in the fitted log outcome:
\[
\Delta\widehat{\log y}=[6+0.02(11)]-[6+0.02(10)]=0.02.
\]
\item Exponentiate to obtain the ratio of the fitted outcomes:
\[
\frac{\tilde y(11)}{\tilde y(10)}=\frac{e^{6.22}}{e^{6.20}}=e^{0.02}=1.020201\ldots
\]
\item Subtract 1 and multiply by 100 to obtain the exact percentage change:
\[
100\frac{\tilde y(11)-\tilde y(10)}{\tilde y(10)}
=100(e^{0.02}-1)=\boxed{2.0201\%}.
\]
\item Approximate using $e^z-1\approx z$ for a small log change $z$:
\[
\%\Delta\tilde y\approx100(0.02\Delta x)=100(0.02)(1)=\boxed{2\%}.
\]
\end{enumerate}
\end{frame}

\begin{frame}{Solution 2: Model C, step by step}
\small
\setlength{\abovedisplayskip}{4pt}
\setlength{\belowdisplayskip}{4pt}
\answerpart{1}C is log--log: 1\% more income is associated with approximately $0.30$\% more in $\tilde y$; its elasticity is $0.30$.
\par\smallskip
\answerpart{2}Use $\widehat{\log y}(x)=6+0.30\log x$.
\begin{enumerate}
\item Subtract the fitted log outcomes; the intercept cancels:
\[
\Delta\widehat{\log y}=0.30(\log11-\log10)=0.30\log(1.10)=0.028593\ldots
\]
\item Exponentiate to obtain the ratio, using $e^{a\log b}=b^a$:
\[
\frac{\tilde y(11)}{\tilde y(10)}=e^{0.30\log(1.10)}=(1.10)^{0.30}=1.029006\ldots
\]
\item Subtract 1 and multiply by 100 for the exact percentage change:
\[
\%\Delta\tilde y=100[(1.10)^{0.30}-1]=\boxed{2.9006\%}.
\]
\item Approximate using $e^z-1\approx z$ and $\log(1+g)\approx g$:
\[
\%\Delta\tilde y\approx100(0.30)\frac{\Delta x}{x}
=100(0.30)(0.10)=\boxed{3\%}.
\]
\end{enumerate}
The last line is elasticity $\times$ percentage change in income: $0.30\times10\%=3\%$.
\end{frame}

\begin{frame}{Solution 2: A doubling is a large change}
\small
\answerpart{3}From 10 to 20, $\Delta x=10$ and $\Delta x/x=1$ (a 100\% increase).
\begin{center}
\begin{tabular}{lll}
\toprule
Model & Exact fitted change & Small-change approximation\\
\midrule
A & $20\log2=13.8629$ points & $20(1)=20$ points\\[4pt]
B & $100(e^{0.20}-1)=22.1403$\% & $100(0.20)=20$\%\\[4pt]
C & $100(2^{0.30}-1)=23.1144$\% & $0.30(100)=30$\%\\
\bottomrule
\end{tabular}
\end{center}
\par\smallskip
For A and C, $\log2=0.6931$ differs substantially from the proportional change 1. For B, exponentiating a log change of $0.20$ gives a proportional change of $0.2214$.
\par\medskip
The log identities are exact; replacing log changes by proportional changes is an approximation for small changes.
\par\smallskip
For B and C, these calculations concern the exponentiated fitted log outcome. It is not generally the conditional arithmetic mean of the original outcome without a retransformation adjustment.
\end{frame}
\fi

\section{Interactions}
\begin{frame}{Problem 3: Does the class-size association differ across districts?}
\small
In 600 schools, let $X$ be class size and $D=1$ identify a low-income district ($D=0$ otherwise). The fitted model is
\[
\widehat{Score}=700-2X-30D+1.20(DX).
\]
Class sizes in both groups range from 15 to 35. The supplied standard error of the interaction estimate is $0.50$. For the joint null that the $D$ and $DX$ coefficients are zero, software reports $F=4.20$ using an appropriate supplied covariance estimate.
\begin{enumerate}
\item Write the fitted line for each group. Interpret the class-size slope in each group and the interaction coefficient.
\item At class size 20, calculate both predicted scores and their difference. Express the low-income-minus-other-district gap as a function of $x$, and find where the two fitted lines cross.
\item State and test the null of equal slopes at 5\%. Then state the null of identical lines, count its restrictions, and use the supplied joint statistic. Explain why these nulls differ.
\end{enumerate}
Use the large-sample critical values $1.96$ for a two-sided $t$ test and $3.00$ for an $F$ test with two restrictions.
\end{frame}

\ifshowsolutions
\begin{frame}{Solution 3: Two group-specific regression lines}
\small
\answerpart{1}Substituting the group indicator gives
\[
\begin{aligned}
D=0:\ &\widehat{Score}=700-2x,\\
D=1:\ &\widehat{Score}=(700-30)+(-2+1.20)x=670-0.80x.
\end{aligned}
\]
An extra student per class is associated with 2 fewer predicted score points for $D=0$ and 0.80 fewer for $D=1$. The interaction $1.20$ is the difference between these slopes: the association is less negative in low-income districts.
\par\medskip
\answerpart{2}At $x=20$, the predictions are $660$ and $654$ points, respectively. The low-income-minus-other gap is $-6$ points.
\[
\widehat{Gap}(x)=(670-0.80x)-(700-2x)=-30+1.20x.
\]
The lines cross at $x=30/1.20=25$, with both predicting 650 points. This lies inside the observed range. The $-30$ intercept difference is the gap at $x=0$, outside that range.
\end{frame}

\begin{frame}{Solution 3: Equal slopes versus identical lines}
\small
Write the model as
\[
Score_i=\beta_0+\beta_X X_i+\beta_D D_i+\beta_{DX}D_iX_i+U_i.
\]
\answerpart{3}Equal slopes require one restriction:
\[
H_0:\beta_{DX}=0,\text{ against }H_1:\beta_{DX}\ne0.
\]
\[
t=\frac{1.20}{0.50}=2.40.
\]
Reject equal slopes at 5\%, since $2.40>1.96$.
\par\medskip
Identical lines require both equal intercepts and equal slopes:
\[
H_0:\beta_D=\beta_{DX}=0.
\]
There are two restrictions. Reject this joint null because the supplied $F=4.20>3.00$.
\par\smallskip
The equal-slope null permits different intercepts. The identical-line null rules out both kinds of difference. These are tests of conditional associations; causal interpretations require additional justification.
\end{frame}
\fi
\end{document}
